\section{Half-chain Schmidt spectrum and the dual fixed point}%
\label{sec:half_chain_schmidt}

This section carries out the Schmidt-decomposition step behind the half-chain
interpretation of the positive diagonal fixed point $\Lambda$ of
\cite[Theorem~6]{PerezGarcia2007Matrix}; it does not establish the full
statement of that theorem. On a ring of $2L$ sites the matrix product vector
splits at the two boundaries into a pair of half-chain families, the overlaps
of the half-chain vectors are matrix elements of the $L$-fold transfer map,
and the unnormalized half-chain reduced density matrix factors exactly through
the half-chain Gram matrix. Away from zero, its spectrum is that of a boundary
comparison matrix of size $D^2$, and under the generic spectral condition of
the source that comparison matrix converges entrywise, geometrically in $L$,
to the diagonal matrix with entries $\lambda_\alpha\lambda_\beta$ up to the
trace normalization of $\Lambda$. Two conclusions of the source remain outside
this section: the passage from entrywise convergence of the comparison matrix
to convergence of its eigenvalues, and the normalization of the finite-size
reduced density matrix to unit trace. The finite-size convention throughout is
the ring of $N=2L$ sites split into two halves of $L$ sites each.

The source states its canonical form with a trace-preserving dual map and a
positive diagonal fixed point $\Lambda$ of that dual. Here the trace-preserving
direction is placed on the transfer map itself and $\Lambda$ is its fixed
point; the two conventions are exchanged by passing to the
conjugate-transposed tensor.

\begin{definition}[Half-chain boundary families]%
    \label{def:half_chain_family}
    \lean{MPSTensor.halfChainMatrix}
    \lean{MPSTensor.halfChainMatrix_apply}
    \lean{MPSTensor.halfChainSwapMatrix}
    \lean{MPSTensor.halfChainSwapMatrix_apply}
    \leanok
    \uses{def:mps_tensor, def:eval_word}
    For a tensor $A$ and a half-chain length $L$, the \emph{half-chain family}
    is the matrix $\Psi$ with rows indexed by configurations
    $\sigma\in\Fin{d}^L$ and columns indexed by boundary pairs
    $(\alpha,\beta)$, whose column at $(\alpha,\beta)$ is the half-chain
    vector
    \begin{align}
        \ket{\Psi_{\alpha,\beta}}
        &=\sum_{\sigma}
        \bra{\alpha}A^{\sigma}\ket{\beta}\,\ket{\sigma}.
        \label{eq:half_chain_vector}
    \end{align}
    The \emph{complementary family} $\Phi$ is the matrix with rows indexed by
    boundary pairs and columns by configurations, whose row at
    $(\alpha,\beta)$ is the vector $\ket{\Psi_{\beta,\alpha}}$ with the two
    boundary indices swapped.
\end{definition}

\begin{definition}[Half-chain coefficient kernel]%
    \label{def:half_chain_kernel}
    \lean{MPSTensor.halfChainKernel}
    \lean{MPSTensor.halfChainKernel_apply}
    \leanok
    \uses{def:mpv}
    The \emph{coefficient kernel} of the ring of $2L$ sites is the matrix over
    the two half-chain configuration spaces whose entry at $(\sigma,\tau)$ is
    the matrix product vector component
    $V^{(2L)}(A)_{\sigma\tau}=\tr(A^{\sigma}A^{\tau})$ of the joined
    configuration.
\end{definition}

\begin{lemma}[Splitting of the ring state at the two boundaries]%
    \label{lem:half_chain_kernel_split}
    \lean{MPSTensor.halfChainKernel_eq_mul}
    \leanok
    \uses{def:half_chain_family, def:half_chain_kernel}
    The coefficient kernel factors as the product $\Psi\Phi$ of the
    half-chain family with its complementary family. Equivalently,
    \begin{align}
        \ket{\psi}
        &=\sum_{\alpha,\beta}
        \ket{\Psi_{\alpha,\beta}}\ket{\Psi_{\beta,\alpha}}.
        \label{eq:half_chain_resolution}
    \end{align}
\end{lemma}

\begin{proof}
    \leanok
    Insert the resolution of the identity
    $\sum_{\beta}\ket{\beta}\bra{\beta}$ at the two cut bonds of the ring
    trace: $\tr(A^{\sigma}A^{\tau})
    =\sum_{\alpha,\beta}\bra{\alpha}A^{\sigma}\ket{\beta}
    \bra{\beta}A^{\tau}\ket{\alpha}$.
\end{proof}

\begin{definition}[Half-chain reduced density matrix]%
    \label{def:half_chain_reduced}
    \lean{MPSTensor.halfChainReducedMatrix}
    \lean{MPSTensor.halfChainReducedMatrix_apply}
    \lean{MPSTensor.posSemidef_halfChainReducedMatrix}
    \leanok
    \uses{def:half_chain_kernel}
    The \emph{unnormalized half-chain reduced density matrix} of the ring
    state on $2L$ sites is $\rho_L=KK^{\dagger}$ with $K$ the coefficient
    kernel. Entrywise this is the partial trace of the unnormalized ring state
    over the complementary half,
    \begin{align}
        (\rho_L)_{\sigma\sigma'}
        &=\sum_{\tau}
        V^{(2L)}(A)_{\sigma\tau}\,
        \overline{V^{(2L)}(A)_{\sigma'\tau}},
        \label{eq:half_chain_reduced_entry}
    \end{align}
    and $\rho_L$ is positive semidefinite.
\end{definition}

\begin{definition}[Half-chain Gram matrix]%
    \label{def:half_chain_gram}
    \lean{MPSTensor.halfChainGram}
    \leanok
    \uses{def:half_chain_family}
    The \emph{half-chain Gram matrix} is $G=\Psi^{\dagger}\Psi$, the matrix of
    overlaps
    $G_{(\alpha,\beta),(\alpha',\beta')}
    =\braket{\Psi_{\alpha,\beta}}{\Psi_{\alpha',\beta'}}$ of the half-chain
    vectors.
\end{definition}

\begin{lemma}[Half-chain overlaps through the transfer map]%
    \label{lem:half_chain_gram_transfer}
    \lean{MPSTensor.halfChainGram_apply_eq_transferMap_pow}
    \leanok
    \uses{def:half_chain_gram, def:transfer_map}
    The half-chain overlaps are matrix elements of the $L$-fold transfer map
    on matrix units:
    \begin{align}
        \braket{\Psi_{\alpha,\beta}}{\Psi_{\alpha',\beta'}}
        &=\bra{\alpha'}\E_A^{L}
        \bigl(\ket{\beta'}\bra{\beta}\bigr)\ket{\alpha}.
        \label{eq:half_chain_gram_transfer}
    \end{align}
\end{lemma}

\begin{proof}
    \leanok
    Expand the overlap as a sum over configurations of products
    $\bra{\alpha'}A^{\sigma}\ket{\beta'}
    \overline{\bra{\alpha}A^{\sigma}\ket{\beta}}$ and recognize the entry of
    $\sum_{\sigma}A^{\sigma}\ket{\beta'}\bra{\beta}(A^{\sigma})^{\dagger}$,
    which is the $L$-fold transfer map applied to the matrix unit.
\end{proof}

\begin{theorem}[Half-chain Schmidt factorization]%
    \label{thm:half_chain_schmidt_factorization}
    \lean{MPSTensor.halfChainReducedMatrix_eq_mul_swap_gram}
    \lean{MPSTensor.halfChainSwapMatrix_mul_conjTranspose}
    \leanok
    \uses{def:half_chain_reduced, def:half_chain_gram}
    Let $S$ denote the swap of the two boundary indices. The unnormalized
    half-chain reduced density matrix factors exactly as
    \begin{align}
        \rho_L
        &=\Psi\,\bigl(SG^{\mathsf T}S\bigr)\,\Psi^{\dagger},
        \label{eq:half_chain_factorization}
    \end{align}
    where $SG^{\mathsf T}S$ is the Gram matrix of the complementary family.
\end{theorem}

\begin{proof}
    \leanok
    \uses{lem:half_chain_kernel_split}
    By the splitting $K=\Psi\Phi$ of
    Lemma~\ref{lem:half_chain_kernel_split},
    $\rho_L=\Psi(\Phi\Phi^{\dagger})\Psi^{\dagger}$, and the Gram matrix
    $\Phi\Phi^{\dagger}$ of the complementary family is the transpose of the
    boundary-swapped half-chain Gram matrix.
\end{proof}

\begin{theorem}[Boundary comparison of the half-chain spectrum]%
    \label{thm:half_chain_charpoly}
    \lean{MPSTensor.charpoly_halfChainReducedMatrix}
    \leanok
    \uses{def:half_chain_reduced, def:half_chain_gram}
    The characteristic polynomials of the half-chain reduced density matrix
    and of the boundary comparison matrix $(SG^{\mathsf T}S)G$ of size $D^2$
    agree after matching the trivial kernel factors:
    \begin{align}
        X^{D^2}\,\chi_{\rho_L}(X)
        &=X^{d^L}\,\chi_{(SG^{\mathsf T}S)G}(X).
        \label{eq:half_chain_charpoly}
    \end{align}
    In particular the nonzero spectrum of $\rho_L$, with multiplicities, is
    that of the boundary comparison matrix.
\end{theorem}

\begin{proof}
    \leanok
    \uses{thm:half_chain_schmidt_factorization,
        lem:matrix_rectangular_charpoly_cyclicity}
    Apply the rectangular cyclicity of the characteristic polynomial to the
    factorization (\ref{eq:half_chain_factorization}), regrouping the product
    as $\Psi$ times $(SG^{\mathsf T}S)\Psi^{\dagger}$; the reversed product is
    $(SG^{\mathsf T}S)\Psi^{\dagger}\Psi=(SG^{\mathsf T}S)G$.
\end{proof}

\begin{theorem}[Geometric transfer error of the half-chain overlaps]%
    \label{thm:half_chain_gram_geometric}
    \lean{MPSTensor.exists_geometric_bound_halfChainGram_sub_fixedPointProj}
    \leanok
    \uses{def:half_chain_gram, def:transfer_map}
    Let $A$ have trace-preserving transfer map $\E_A$ with fixed point
    $\Lambda$, $\tr(\Lambda)\neq0$, and suppose
    \begin{align}
        \rho_{\mathrm{spec}}\bigl(\E_A-P_{\Lambda}\bigr)&<1,
        \label{eq:half_chain_c2_gap}
    \end{align}
    where $P_{\Lambda}$ is the fixed-point projection. Then there are
    constants $C>0$ and $0<r<1$ such that for every $L\geq1$ and all boundary
    pairs,
    \begin{align}
        \Bigl|
        G_{(\alpha,\beta),(\alpha',\beta')}
        -\bra{\alpha'}P_{\Lambda}
        \bigl(\ket{\beta'}\bra{\beta}\bigr)\ket{\alpha}
        \Bigr|
        &\leq C r^{L}.
        \label{eq:half_chain_gram_geometric}
    \end{align}
    The hypothesis (\ref{eq:half_chain_c2_gap}) expresses the generic
    condition C2 of \cite{PerezGarcia2007Matrix} in the trace-preserving
    gauge: apart from the simple eigenvalue one, the spectrum of $\E_A$ lies
    strictly inside the unit disc, and the correction term is of the order of
    the subleading spectral radius, the bound of order $|\nu_2|^{N/2}$ of the
    source.
\end{theorem}

\begin{proof}
    \leanok
    \uses{lem:half_chain_gram_transfer,
        thm:geometric_bound_of_spectral_radius_lt_one}
    Write $\E_A^{L}=P_{\Lambda}+(\E_A-P_{\Lambda})^{L}$ for $L\geq1$ by the
    power decomposition. By (\ref{eq:half_chain_gram_transfer}) the deviation
    of a Gram entry from the fixed-point projection entry is an entry of
    $(\E_A-P_{\Lambda})^{L}$ applied to a matrix unit, bounded by the operator
    norm of the $L$-th power, which decays geometrically below unit spectral
    radius.
\end{proof}

\begin{lemma}[Diagonal limit of the half-chain overlaps]%
    \label{lem:half_chain_gram_limit_diagonal}
    \lean{MPSTensor.fixedPointProj_single_entry_diagonal}
    \leanok
    For a diagonal fixed point
    $\Lambda=\operatorname{diag}(\lambda_1,\ldots,\lambda_D)$ with
    $\tr(\Lambda)\neq0$, the limiting overlaps form a diagonal matrix in the
    boundary pairs:
    \begin{align}
        \bra{\alpha'}P_{\Lambda}
        \bigl(\ket{\beta'}\bra{\beta}\bigr)\ket{\alpha}
        &=\frac{\lambda_\alpha}{\tr(\Lambda)}\,
        \delta_{\alpha,\alpha'}\delta_{\beta,\beta'}.
        \label{eq:half_chain_limit_diagonal}
    \end{align}
    For $\tr(\Lambda)=1$ this is the display
    $\lambda_\alpha\delta_{\alpha,\alpha'}\delta_{\beta,\beta'}$ of
    \cite{PerezGarcia2007Matrix}.
\end{lemma}

\begin{proof}
    \leanok
    Evaluate the fixed-point projection on the matrix unit: its trace is
    $\delta_{\beta,\beta'}$, and the diagonal entry of $\Lambda$ at
    $(\alpha',\alpha)$ is $\lambda_\alpha\delta_{\alpha,\alpha'}$.
\end{proof}

\begin{theorem}[Geometric convergence to the doubled fixed-point spectrum]%
    \label{thm:half_chain_comparison_geometric}
    \lean{MPSTensor.exists_geometric_bound_halfChainComparison_sub_diagonal}
    \leanok
    \uses{def:half_chain_gram, def:transfer_map}
    Let $A$ have trace-preserving transfer map with diagonal fixed point
    $\Lambda=\operatorname{diag}(\lambda_1,\ldots,\lambda_D)$,
    $\tr(\Lambda)\neq0$, satisfying the spectral condition
    (\ref{eq:half_chain_c2_gap}). Then there are constants $C>0$ and $0<r<1$
    such that for every $L\geq1$ the boundary comparison matrix satisfies,
    entrywise,
    \begin{align}
        \Bigl|
        \bigl((SG^{\mathsf T}S)G\bigr)_{(\alpha,\beta),(\alpha',\beta')}
        -\frac{\lambda_\alpha\lambda_\beta}{\tr(\Lambda)^2}\,
        \delta_{\alpha,\alpha'}\delta_{\beta,\beta'}
        \Bigr|
        &\leq C r^{L}.
        \label{eq:half_chain_comparison_geometric}
    \end{align}
    The limiting matrix is the diagonal matrix $\Lambda\otimes\Lambda$ up to
    the trace normalization. Together with
    Theorem~\ref{thm:half_chain_charpoly}, the nonzero spectrum of the
    half-chain reduced density matrix on the ring of $2L$ sites is that of a
    $D^2\times D^2$ matrix converging entrywise, geometrically in $L$, to
    $\Lambda\otimes\Lambda/\tr(\Lambda)^2$. The remaining passage to the
    convergence of the eigenvalues themselves is the continuity of the
    spectrum of a matrix in its entries and is not part of this section.
\end{theorem}

\begin{proof}
    \leanok
    \uses{thm:half_chain_gram_geometric, lem:half_chain_gram_limit_diagonal}
    Both factors of the comparison matrix are entrywise within $Cr^{L}$ of
    their diagonal limits by Theorem~\ref{thm:half_chain_gram_geometric} and
    Lemma~\ref{lem:half_chain_gram_limit_diagonal}, the swapped-transposed
    factor because swapping and transposing permute entries and exchange the
    two diagonal weights. The limit entries are bounded, so the product
    deviates from the product of the limits by at most a constant multiple of
    $r^{L}$, with constant $D^2\bigl(C(B+C)+BC\bigr)+1$ for an entry bound
    $B$ on the limits.
\end{proof}
